%An example is shown at \figref{dataflow-feature-by-feature-correction};

\begin{figure}
    \begin{tikzpicture}[every node/.style={font=\sffamily, scale=0.6, text width=7cm}, align=left]
        \node (start) [code, label=below:Initial feature in class] {
            \begin{lstlisting}
class
    EXAMPLE
feature
    min (x, y: INTEGER): INTEGER
        -- Header comment
        do
            Result := x
        ensure
            res: (Result = x or Result = y)
                and Result <= x and Result <= y
        end

    max (x, y: INTEGER): INTEGER
        -- Header comment
        do
            Result := x
        ensure
            res: (Result = x or Result = y)
                and Result >= x and Result >= y
        end
end
\end{lstlisting}
        };

        \node (prompt) [code, right=of start, yshift=-190pt,label=below:Generated prompt] {
\begin{lstlisting}
class
    EXAMPLE
feature
    -- <task>
    min (x, y: INTEGER): INTEGER
        -- Header comment
        do
            Result := x
        ensure
            res: (Result = x or Result = y)
                and Result <= x and Result <= y
        end
        -- AutoProof counterexample

    -- <task>
    max (x, y: INTEGER): INTEGER
        -- Header comment
        do
            Result := x
        ensure
            res: (Result = x or Result = y)
                and Result >= x and Result >= y
        end
        -- AutoProof counterexample
end
\end{lstlisting}
        };
        \draw[->, thick] (start) -- (prompt);

        \node (llm-candidate) [code, left=of prompt, yshift=-140pt, label=below:Generated llm-candidate] {
            \begin{lstlisting}
class
    EXAMPLE
feature
    -- <task>
    min (x, y: INTEGER): INTEGER
        -- Header comment
        do
            if x < y then
                Result := x
            else
                Result := y
            end
        ensure
            res: (Result = x or Result = y)
                and Result <= x and Result <= y
        end
        -- AutoProof counterexample

    -- <task>
    max (x, y: INTEGER): INTEGER
        -- Header comment
        do
            if x < y then
                Result := y
            else
                Result := x
            end
        ensure
            res: (Result = x or Result = y)
                and Result >= x and Result >= y
        end
        -- AutoProof counterexample
end
\end{lstlisting}
        };
        \draw[->, thick] (prompt) -- (llm-candidate);

        \node (verification-result) [code, right=of llm-candidate, yshift=-150pt] { Verification output };
        \draw[->, thick] (llm-candidate) -- (verification-result);
        \draw[->, thick] (verification-result.east) -| ([xshift=30pt] prompt.east)
            node[midway, below, yshift=100pt] {Reset failing features}
            node[near start, above, xshift=60pt] {Some feature fails} -- (prompt.east);
        \node (initial-feature-after-edits) [code, left=of verification-result, yshift=-240pt, label=below:Feature after edits] {
            \begin{lstlisting}
class
    EXAMPLE
feature
    min (x, y: INTEGER): INTEGER
        -- Header comment
        do
            if x < y then
                Result := x
            else
                Result := y
            end
        ensure
            res: (Result = x or Result = y)
                and Result <= x and Result <= y
        end

    max (x, y: INTEGER): INTEGER
        -- Header comment
        do
            if x < y then
                Result := y
            else
                Result := x
            end
        ensure
            res: (Result = x or Result = y)
                and Result >= x and Result >= y
        end
end
\end{lstlisting}
        };
        \draw[->, thick] (verification-result) -- (initial-feature-after-edits)
            node[midway, right] {All features verify};
    \end{tikzpicture}
    \caption{Dataflow of a feature in the feature by feature correction command}\label{dataflow-feature-by-feature-correction}
\end{figure}
